% =====================================================================
% Section to insert into Paper 7(a) — Bound Core Electron Model.
% Second draft, 2026-06-12 (first draft in git history).
% Recommended placement: as a new \section{} between §10 (Coulomb's
% Constant Derivation) and the existing §11 (Open Questions). Renumber
% Open Questions to §12.
%
% Revision notes (vs. first draft):
%   (1) The two windings are distinguished: azimuthal pattern advance at
%       omega_C (phase motion, carries the charge cycle -> I_eff, mu_B)
%       vs. poloidal material rotation at Omega = m_A c^2/hbar (inherited
%       from the host vortex flow -> the Coulomb source clock).
%   (2) The mu_0 paragraph is rewired: the deferral to the longitudinal
%       A = xi amplitude closure is superseded by the winding-pump
%       normalization (q_0 = 2 pi r_loop), which is the SAME
%       one-winding-per-cycle postulate used for I_eff here.
%   (3) New subsection on species scaling: mu scales as 1/m across the
%       leptons (observed) precisely because e is a universal input —
%       stated as a structural constraint on the charge-source derivation.
%
% NOTES FOR INSERTION:
%   - Cross-references \ref{sec:force-constants},
%     \ref{sec:form-factor-topological}, \ref{sec:qss} are placeholders —
%     verify against the actual labels in 7(a) (§10 is
%     \label{sec:coulomb-derivation}) and Paper 7.
%   - The quoted Coulomb accuracy (-0.07%) refers to the spin-channel
%     (winding-pump) formulation at the joint parameter point
%     (rho_A ~ 5e11 kg/m^3, m_A ~ 666 MeV/c^2). If this section is
%     inserted before §10 is revised to that formulation, quote the
%     current §10 figure (4%) instead.
%   - r_{loop} is written out as in the first draft; harmonize to the
%     \rloop macro when inserting.
% =====================================================================

\section{Structural Confirmation: The Bohr Magneton from Ring Geometry}
\label{sec:bohr-magneton}

The derivation of §\ref{sec:force-constants} fixes the Coulomb constant
$k$ from the bound core ring's spin-phase source amplitude. A
complementary, parameter-free check on the ring geometry comes from
computing its classical magnetic moment and comparing to the observed
Bohr magneton $\mu_B = e\hbar/(2m_e)$. The match is exact: no fit,
no closure assumptions, no source-identification factor.

\subsection{Effective Current of the Bound Core Ring}

The ring carries charge $e$ as a topological winding number
(§\ref{sec:form-factor-topological}). The winding \emph{pattern}
advances azimuthally around the loop once per Compton period
$T_C = 2\pi\hbar/(m_e c^2)$. This advance is a phase motion, not a
material flow: the pattern speed at the loop radius is
$\omega_C\,r_{\text{loop}} = c$ exactly, while the host circulation at
that radius moves constituents at only
$v(r_{\text{loop}}) = m_e c/(2 m_A) \approx 4\times10^{-4}\,c$. (Were
the advance material transport, the combined azimuthal and poloidal
speeds of the core constituents would exceed $c$; the winding advance
is the motion of a pattern in the relative phase, not of the
constituents that carry it.) An external observer averaging over the
loop therefore sees one quantum of charge pass any azimuthal
cross-section per Compton period:
\begin{equation}
I_{\text{eff}} = \frac{e}{T_C}
              = \frac{e\,\omega_C}{2\pi}
              = \frac{e\,m_e c^2}{2\pi\,\hbar}.
\label{eq:I-eff-bohr}
\end{equation}

\subsection{Two Windings, Two Clocks}
\label{sec:two-clocks}

The azimuthal advance above is one of \emph{two} rotations the bound
structure carries, and they must not be conflated. The same
relative-phase winding also rotates poloidally---about the core tube's
own axis---at the rate inherited from the host vortex flow at the
formation boundary:
\begin{equation}
\Omega \;=\; \frac{v(\xi)}{\xi}
        \;=\; \frac{c/\sqrt{2}}{\hbar/(\sqrt{2}\,m_A c)}
        \;=\; \frac{m_A c^2}{\hbar},
\label{eq:poloidal-clock}
\end{equation}
exact because $v(\xi) = c/\sqrt{2}$ is itself exact and
$m_A$-independent in construction: $\Omega$ is a property of the
medium, common to every singly-wound structure regardless of its loop
radius. The two clocks differ by
$\Omega/\omega_C = m_A/m_e \approx 1.3\times10^{3}$: the winding
re-threads poloidally some thirteen hundred times per azimuthal
cycle---a fine-pitch screw. The division of labor is clean. The
poloidal rotation is material (the bound core co-rotates with the flow
that created it, at $v(\xi) = c/\sqrt{2}$, safely subluminal) and is
the clock of the Coulomb source (§\ref{sec:force-constants}): its
uniform phase precession at the boundary drives the spin-phase pump.
The azimuthal advance is pattern motion and is the clock of the
magnetic moment computed here. The poloidal rotation is also
compatible with everything below: its local rotation axes point
azimuthally and close around the loop, so it contributes no net
angular momentum about the symmetry axis and leaves both $\mu_M$ and
$S$ untouched.

\subsection{Magnetic Moment of the Loop}

Treating the ring as a closed current loop of radius $r_{\text{loop}}$
and area $\mathcal{A} = \pi r_{\text{loop}}^2$, the classical magnetic
moment is $\mu_M = I_{\text{eff}}\,\mathcal{A}$. Substituting
$r_{\text{loop}} = \hbar/(m_e c)$ from Paper 7:
\begin{equation}
\mu_M = \frac{e\,m_e c^2}{2\pi\,\hbar}\cdot \pi\cdot\frac{\hbar^2}{m_e^2 c^2}
      = \frac{e\,\hbar}{2 m_e}
      = \mu_B.
\label{eq:bohr-magneton}
\end{equation}
The cancellation in~(\ref{eq:bohr-magneton}) is exact and parameter-free.

\subsection{What the Match Confirms}

Equation~(\ref{eq:bohr-magneton}) requires three independently-motivated
structural inputs to combine in precisely the right way:
\begin{itemize}
\item The ring radius is the reduced Compton wavelength:
  $r_{\text{loop}} = \hbar/(m_e c)$. Any other length scale produces the
  wrong answer.
\item The azimuthal clock is the Compton frequency:
  $\omega_C = m_e c^2/\hbar$. This is the frequency of the winding
  pattern's advance around the loop (Paper 7), distinct from the
  poloidal clock $\Omega$ of §\ref{sec:two-clocks}.
\item The effective current is one charge per period:
  $I_{\text{eff}} = e\,\omega_C/(2\pi)$. This follows from the
  topological character of charge established in
  §\ref{sec:form-factor-topological}: the ring carries a single winding
  whose azimuthal pattern advances exactly once per cycle.
\end{itemize}
Each input is motivated independently elsewhere in the framework, and
their product yielding exactly $\mu_B$ is a non-trivial geometric
self-consistency. As developed below, the third input is not unique to
this calculation: the same one-winding-per-cycle normalization
independently fixes the Coulomb source amplitude, making the pair of
results a single-postulate structure rather than a coincidence.

\subsection{Gyromagnetic Ratio $g = 2$}

The same construction gives the Dirac value of the gyromagnetic ratio
without further assumptions. The bound core ring carries
quantized half-circulation $\kappa_e = h/(2 m_A)$
(cf.\ Paper~7, §\ref{sec:qss}), which corresponds to spin angular
momentum $S = \hbar/2$. Combining with~(\ref{eq:bohr-magneton}):
\begin{equation}
\mu_B \;=\; g\cdot\frac{e}{2 m_e}\cdot S \quad\Longrightarrow\quad
g \;=\; 2,
\end{equation}
to leading order. (The identification is exact per constituent---each
participant in the half-quantum flow carries
$m_A \kappa_e/2\pi = \hbar/2$---while the bookkeeping that assigns the
\emph{structure} a total of exactly $\hbar/2$ is carried by Paper 7's
quantized-circulation argument and is not re-derived here.) The
standard QED anomalous corrections
$g = 2 + \alpha/\pi + \mathcal O(\alpha^2)$ are expected to arise from
coupled fluctuations of the surrounding Aether condensate (the analog of
the radiative diagrams in QED) and lie beyond the classical ring
picture.

\subsection{What This Does \emph{Not} Determine}

The match in~(\ref{eq:bohr-magneton}) yields $\mu_B$ in SI units of
$\mathrm{A\cdot m^2}$ from purely geometric and kinematic ring
properties. It does \emph{not} fix the vacuum permeability $\mu_0$.
That constant enters only when one computes the magnetic field
produced by the moment at distance $r$
($B = \mu_0\mu_M/(4\pi r^3)$) or the interaction energy of two such
moments, and it is fixed ($\mu_0 = 4\pi k/c^2$) by the
Coulomb-constant derivation of §\ref{sec:force-constants}. An earlier
draft of this section deferred that derivation to the longitudinal
amplitude closure $A = \xi$; that closure is superseded. The Coulomb
source is the ring's spin-phase winding pump, and its amplitude is
normalized by the \emph{same} postulate used in
Eq.~(\ref{eq:I-eff-bohr}): one winding advance per Compton cycle,
giving source strength $q_0 = 2\pi\,r_{\text{loop}}$---one winding
quantum across the ring aperture per cycle. The two calculations
therefore share their single dimensionless input:
\begin{itemize}
\item one winding per cycle, counted as an azimuthal current
  $\;\rightarrow\;$ $\mu_B$, exactly;
\item one winding per cycle, counted as a spin-phase pump
  $\;\rightarrow\;$ $k$, to $-0.07\%$ at the joint parameter point
  ($\rho_A \approx 5\times10^{11}$~kg/m$^3$,
  $m_A \approx 666$~MeV/$c^2$).
\end{itemize}
One topological postulate, two observables. This is the sense in which
the product in the preceding subsection is quantized rather than
tuned: a normalization adjusted by hand to fit $k$ would have no reason
to land the magneton exactly, and vice versa.

\subsection{Species Scaling: Why $\mu$ Scales and $e$ Must Not}
\label{sec:species-scaling}

The construction extends across the charged leptons by
$r_{\text{loop}} \to \hbar/(m c)$ and $T_C \to 2\pi\hbar/(m c^2)$,
predicting $\mu = e\hbar/(2m)$---as observed: the muon's moment is
$e\hbar/(2m_\mu)$ with the same $g \approx 2$. The scaling succeeds
precisely because $e$ enters~(\ref{eq:I-eff-bohr}) as a universal input
while only the geometry scales with the species mass. This is a
structural constraint on the companion Coulomb derivation: the charge
quantum must \emph{not} inherit the $r_{\text{loop}}$ scaling that the
moment correctly does---a source strength proportional to each
species' own loop radius would make the muon's charge smaller than the
electron's by $m_e/m_\mu$, against observation. The resolution
direction is available within the framework: $m_e$ is itself fixed by
the medium through the self-consistent mass equation
(§\ref{sec:combined}), so $q_0 = 2\pi\hbar/(m_e c)$ is a constant of
the medium---the charge quantum of the condensate's ground ring
mode---to be inherited by any charged structure, whatever its own size
and clocks. Making that inheritance mechanism explicit is an open
problem (§\ref{sec:open-questions}).

\subsection{Status}

\textbf{Derived (D).} Given (i) the ring radius from Paper 7,
(ii) the topological character of charge
(§\ref{sec:form-factor-topological}), and (iii) the half-quantum
circulation from Paper 7 §\ref{sec:qss}, the Bohr magneton and $g = 2$
follow with no free parameters or fits. The shared
one-winding-per-cycle normalization ties this result to the Coulomb
source amplitude, elevating the two to a single-postulate pair.
\textbf{Open:} the total angular-momentum bookkeeping behind
$S = \hbar/2$, and the species-inheritance mechanism for the charge
quantum (§\ref{sec:species-scaling}).
